\section{Quality Metrics}\label{sec:5}

This section defines the scores computed inside the Matchmaking Agent and the ranking, explanation, and fairness criteria used in offline evaluation.
Implementation modules appear in Section~\ref{sec:4}; numeric outcomes appear in Section~\ref{sec:6}.

\subsection{Evaluation corpus and protocol}\label{sec:5-corpus}

Offline evaluation treats each resume as a query and each job posting as a candidate document.
The labeled set comprises 30 resumes, 15 jobs, and 47 graded query--document pairs in \texttt{data/eval\_pairs.json}.
Relevance grades are $r \in \{0,1,2\}$ (none, partial, strong); queries with $r>0$ form the relevant set for binary metrics, while nDCG uses the full graded scale.

The relevance labels were assigned manually by the project authors during demo-corpus construction; they are not inferred from portal clicks or hiring outcomes.
Each candidate--job pair was judged against the structured snapshot fields used by the matcher: resume skills, experience level, job title, required skills, compensation constraints, and remote-policy fit.
A strong label indicates that the candidate satisfies the core role requirements; a partial label indicates plausible fit with missing or weaker requirements; and a none label indicates no clear role fit.
Final labels were chosen conservatively: when a pair was borderline, it was kept at partial or none rather than promoted to strong.

Unless a driver overrides \texttt{top\_k}, metrics are macro-averaged at $K{=}5$ over all 30 queries.
All benchmark drivers use the same metric routine so method comparisons differ only in the ranking function.
Table~\ref{tab:dataset} summarizes corpus statistics.

\begin{center}
\refstepcounter{table}\label{tab:dataset}
\par\smallskip\noindent\textbf{Table~\thetable:} Evaluation corpus statistics.
\par\smallskip
\begin{tabular}{@{}ll@{}}
\toprule
Statistic & Value \\
\midrule
Task & Résumé $\rightarrow$ jobs (candidate-initiated retrieval) \\
Candidates (queries) & 30 \\
Job postings (corpus) & 15 \\
Graded query--document pairs & 47 \\
Queries with $\geq 1$ relevant job & 30 \\
Relevance scale & 0 (none), 1 (partial), 2 (strong) \\
Grade-2 (strong) labels & 21 \\
Grade-1 (partial) labels & 26 \\
Labels per query (min / max) & 1 / 2 \\
Embedding model & \texttt{all-MiniLM-L6-v2} ($d{=}384$) \\
Label source & \path{data/eval_pairs.json} v1.0 \\
\bottomrule
\end{tabular}
\par\vspace{0.4em}
\footnotesize{Starter labeled set for offline benchmarks; not derived from live portal traffic.}
\end{center}

\subsection{Matching scores}\label{sec:5-matching-scores}

A resume and a job snapshot pass through embedding similarity, skill overlap, and the experience, compensation, and remote signals, which combine into a weighted composite score and a structured explanation that drives the ranked recommendation.
All component scores are normalized to $[0,1]$.
Let $\mathbf{e}_c, \mathbf{e}_j \in \mathbb{R}^{d}$ denote candidate and job snapshot embeddings ($d{=}384$), and let $A$ and $B$ be canonical skill sets on the resume and required job skills, respectively.

The semantic score compares resume and job-description embeddings.
Cosine mode is the portal default:
\begin{equation}
  s_{\mathrm{sem}}^{\mathrm{cos}} =
    \frac{\mathbf{e}_c \cdot \mathbf{e}_j}{\|\mathbf{e}_c\|\,\|\mathbf{e}_j\|}.
  \label{eq:semantic-cosine}
\end{equation}
Euclidean mode maps distance to similarity:
\begin{equation}
  s_{\mathrm{sem}}^{\mathrm{euc}} = \frac{1}{1 + \|\mathbf{e}_c - \mathbf{e}_j\|_2}.
  \label{eq:semantic-euclidean}
\end{equation}

The skill score compares resume skills against job requirements.
Jaccard mode is the portal default:
\begin{equation}
  s_{\mathrm{skl}}^{\mathrm{jac}} =
    \begin{cases}
      |A \cap B| / |A \cup B| & \text{if } A \neq \emptyset \text{ and } B \neq \emptyset,\\
      0 & \text{otherwise}.
    \end{cases}
  \label{eq:jaccard}
\end{equation}
Embedding mode averages, for each required skill $b \in B$, the maximum cosine between $b$'s skill phrase embedding and any resume skill embedding:
\begin{equation}
  s_{\mathrm{skl}}^{\mathrm{emb}} =
    \frac{1}{|B|} \sum_{b \in B}
    \max_{a \in A} \cos\!\bigl(\mathbf{e}_a^{(\mathrm{skill})}, \mathbf{e}_b^{(\mathrm{skill})}\bigr).
  \label{eq:skill-embed}
\end{equation}
A third, \emph{graded relation-aware} mode addresses the all-or-nothing behavior of Jaccard by giving partial (never full) credit to related skills. It resolves each required skill $b \in B$ to a relation class against the resume skills and averages the best graded credit over requirements:
\begin{equation}
  s_{\mathrm{skl}}^{\mathrm{grd}} = \frac{1}{|B|}\sum_{b \in B} \max_{a \in A} g(a,b),\quad
  g(a,b) = \begin{cases} 1.0 & \text{same canonical skill (\textsc{exact})},\\ 0.5 & \text{same taxonomy group (\textsc{related})},\\ 0 & \text{otherwise}. \end{cases}
  \label{eq:skill-graded}
\end{equation}
The credits are fixed a priori (not tuned). Holding all six weights fixed, we evaluate this matcher as a drop-in for the skill channel and decompose its effect into the coverage-form change (symmetric Jaccard $\to$ asymmetric required-coverage) and the relation-aware partial credit. On the high-power synthetic corpus the coverage form drives the gain (composite nDCG@5 $0.917 \to 0.949$) while the relation-aware credit does not help there --- the synthetic ground truth is exact coverage and has no related-skill structure. On the real human corpus the coverage form is null but the relation-aware credit improves six of thirty queries and worsens none (sign-test $p{=}0.03$, one query dominant, effective $n{=}6$), which we read as a directional signal on a near-ceiling benchmark rather than a broad effect. The portal retains Jaccard for comparability; the graded matcher is reported as an auditable enhancement whose measurable benefit is specific to human-judged relevance.

A two-signal blend combines semantic and skill scores with weight $\alpha$ (default $\alpha{=}0.7$):
\begin{equation}
  s_{\mathrm{mm}} = \alpha\, s_{\mathrm{sem}} + (1-\alpha)\, s_{\mathrm{skl}}.
  \label{eq:multimodal}
\end{equation}

The composite score is a fixed weighted sum over six interpretable components---semantic, skills, title token overlap, experience fit, compensation fit, and remote preference---then clipped to $[0,1]$:
\begin{equation}
  \begin{aligned}
  s_{\mathrm{final}} &=
    \mathrm{clip}_{[0,1]}\!\left(\sum_{k} w_k s_k\right),\\
  (w_{\mathrm{sem}}, w_{\mathrm{skl}}, w_{\mathrm{tit}}, w_{\mathrm{exp}},
   w_{\mathrm{comp}}, w_{\mathrm{rem}})
    &= (0.28, 0.27, 0.10, 0.15, 0.10, 0.10).
  \end{aligned}
  \label{eq:composite}
\end{equation}
Reciprocal rank fusion (RRF), used in comparison experiments, merges $m$ ranked lists without retraining:
\begin{equation}
  \mathrm{RRF}(x) = \sum_{i=1}^{m} \sum_{r} \mathbb{1}[\mathcal{L}_i[r] = x] \cdot \frac{1}{k + r},
  \label{eq:rrf}
\end{equation}
with smoothing constant $k{=}60$ by default.
Optional constraints, Platt calibration, learned fusion, feedback boosts, cross-encoder reranking, and ANN shortlist sweeps are evaluated as experimental branches; portal defaults keep them disabled.
A held-out calibration comparison finds that Platt scaling reaches expected calibration error (ECE) 0.019 but with low discrimination (Brier skill near zero), whereas beta calibration attains a lower ECE (0.009) under both equal-width and equal-mass binnings \emph{and} preserves discrimination (Brier skill 0.67); we report this trade-off honestly and recommend beta as the confidence head while retaining the simpler Platt map as the frozen default.

\subsection{Ranking metrics}\label{sec:5-ranking}

Given ranked list $\pi_q$ for query $q$, graded relevance map $\mathrm{rel}_q(\cdot)$, and relevant set $\mathcal{R}_q = \{ j : \mathrm{rel}_q(j) > 0 \}$, let $\pi_q[:K]$ denote the top-$K$ prefix.

Precision@K is the fraction of the top-$K$ slots occupied by relevant jobs:
\begin{equation}
  \mathrm{P@}K(q) = \frac{1}{K}\sum_{j \in \pi_q[:K]} \mathbb{1}[j \in \mathcal{R}_q].
  \label{eq:p-at-k}
\end{equation}

Recall@K is the fraction of all relevant jobs for the query retrieved within the top-$K$:
\begin{equation}
  \mathrm{R@}K(q) = \frac{1}{|\mathcal{R}_q|}\sum_{j \in \pi_q[:K]} \mathbb{1}[j \in \mathcal{R}_q].
  \label{eq:r-at-k}
\end{equation}

Normalized discounted cumulative gain rewards higher grades at higher ranks.
For graded relevance:
\begin{equation}
  \mathrm{DCG@}K(q) = \sum_{i=1}^{K} \frac{2^{\mathrm{rel}_q(j_i)} - 1}{\log_2(i+1)},
  \label{eq:dcg}
\end{equation}
\begin{equation}
  \mathrm{nDCG@}K(q) = \frac{\mathrm{DCG@}K(q)}{\mathrm{DCG@}K(q^\ast)},
  \label{eq:ndcg}
\end{equation}
where $\pi_{q^\ast}$ is the ideal ranking that sorts jobs by $\mathrm{rel}_q$ descending.

MRR treats relevance as binary ($r>0$) and measures how early the first relevant job appears:
\begin{equation}
  \mathrm{RR}(q) =
    \begin{cases}
      1 / \mathrm{rank}_q(j^\ast) & \text{if a relevant job } j^\ast \text{ appears at rank } \mathrm{rank}_q(j^\ast),\\
      0 & \text{otherwise}.
    \end{cases}
  \label{eq:mrr}
\end{equation}
Macro-averaged $\mathrm{MRR} = \frac{1}{|\mathcal{Q}|}\sum_{q \in \mathcal{Q}} \mathrm{RR}(q)$.

MAP averages precision at each relevant hit, normalized by $|\mathcal{R}_q|$:
\begin{equation}
  \mathrm{AP}(q) = \frac{1}{|\mathcal{R}_q|}
    \sum_{r=1}^{|\pi_q|}
    \mathrm{P@}r(q) \cdot \mathbb{1}[\pi_q[r] \in \mathcal{R}_q],
  \label{eq:map}
\end{equation}
where $\mathrm{P@}r(q)$ is precision computed over the length-$r$ prefix.
Drivers report macro-averaged P@K, R@K, MRR, nDCG@K, and MAP unless noted otherwise.

\subsection{Benchmark comparison criteria}\label{sec:5-benchmark-criteria}

The benchmark suite evaluates, on the \emph{same} labeled corpus: lexical baselines, dense strategies, composite and component ablations, fusion variants, optional cross-encoder reranking, and ANN shortlist sweeps.
Primary comparison metrics are macro-averaged P@5, R@5, MRR, nDCG@5, and MAP at $K{=}5$.

Paired bootstrap tests resample queries with replacement ($N{=}5000$, seed~42) and compare method pairs on nDCG@K and MRR; significant comparisons are reported in Section~\ref{sec:6-main-result}.

Continuous integration recomputes selected rankings and requires:
\begin{equation}
  \left| \widehat{\mathrm{nDCG@}5}_{\mathrm{run}} - \mathrm{nDCG@}5_{\mathrm{expected}} \right| \leq 0.04,
  \label{eq:regression}
\end{equation}
against committed values in \path{data/expected/paper_progression_summary.json} (\path{tests/benchmarks/test_eval_regression.py}).

From the top-$P{=}10$ multimodal pool per query, hard negatives are high-scoring jobs outside the labeled relevant set:
\begin{equation}
  \mathcal{H}_q = \left\{ j \in \pi_q[:P] : j \notin \mathcal{R}_q \right\}_{|\cdot|\leq 5}.
  \label{eq:hard-negs}
\end{equation}
Success criterion: zero overlap between $\mathcal{H}_q$ and declared relevant jobs across all queries (label-set sanity, not ranking quality).

ANN shortlist latency is summarized in Section~\ref{sec:6-strengths}.

\subsection{Explanation quality criteria}\label{sec:5-explainability}

Explainability is evaluated offline on ranked pairs by comparing rule-based bullets (portal default) with a grounded template explainer.
Faithfulness requires explanations to cite available matched or missing skills, avoid hallucinated skills, and align qualitative claims with the underlying score breakdown.
Specificity measures whether bullets contain concrete skill references rather than generic text.
For synthetic profile variants, consistency is measured with bullet-set Jaccard similarity and matched-skill equality.

Aggregate explanation-quality results are summarized in Section~\ref{sec:6-limitations}.

\subsection{Fairness and sanity metrics}\label{sec:5-fairness}

Fairness checks are proxy-based sanity audits on the synthetic demo corpus; they do not infer protected attributes from real users.
Queries are grouped by experience tier and remote preference to compute selection-rate disparate-impact ratios.
A separate counterfactual audit applies controlled edits to ten fabricated profile pairs and records rank shifts, score deltas, and explanation drift.
Numeric outcomes appear in Section~\ref{sec:6-limitations}.
